% !TeX encoding = UTF-8
\documentclass[14pt,a4paper]{extarticle}

\newif\ifusekprj
\IfFileExists{kprj.sty}{\usekprjtrue\usepackage{kprj}}{\usekprjfalse}
\usepackage{tabularx}
\usepackage{float}
\usepackage{flafter}
\usepackage{graphicx}
\usepackage{booktabs}
\usepackage{longtable}
\usepackage{multirow}
\usepackage{array}
\usepackage{amsmath,amssymb}
\usepackage{iftex}
\ifPDFTeX
  \usepackage[T2A]{fontenc}
  \usepackage[utf8]{inputenc}
  \usepackage[russian]{babel}
\else
  \usepackage{fontspec}
  \usepackage{polyglossia}
  \setdefaultlanguage{russian}
  \setmainfont{CMU Serif}
  \setsansfont{CMU Sans Serif}
  \setmonofont{CMU Typewriter Text}
\fi
\usepackage{hyperref}
\usepackage{listings}
\usepackage{xcolor}

\makeatletter
\ifusekprj
\else
  \newcommand{\authorfull}[1]{\gdef\@authorfull{#1}}
  \newcommand{\faculty}[1]{\gdef\@faculty{#1}}
  \newcommand{\facultyshort}[1]{\gdef\@facultyshort{#1}}
  \newcommand{\chair}[1]{\gdef\@chair{#1}}
  \newcommand{\chairno}[1]{\gdef\@chairno{#1}}
  \newcommand{\group}[1]{\gdef\@group{#1}}
  \newcommand{\chief}[1]{\gdef\@chief{#1}}
  \newcommand{\FirstConsultant}[1]{\gdef\@FirstConsultant{#1}}
  \newcommand{\inspector}[1]{\gdef\@inspector{#1}}
  \newcommand{\workyear}[1]{\gdef\@workyear{#1}}
  \newcommand{\setcontentsname}[1]{\renewcommand{\contentsname}{#1}}
  \newcommand{\setrefsname}[1]{\gdef\@refsname{#1}}
  \providecommand{\@authorfull}{}
  \providecommand{\@faculty}{}
  \providecommand{\@facultyshort}{}
  \providecommand{\@chair}{}
  \providecommand{\@chairno}{}
  \providecommand{\@group}{}
  \providecommand{\@chief}{}
  \providecommand{\@FirstConsultant}{}
  \providecommand{\@inspector}{}
  \providecommand{\@workyear}{\the\year}
  \providecommand{\@refsname}{Список использованных источников}
\fi
\makeatother

\hypersetup{
  colorlinks=true,
  linkcolor=black,
  citecolor=black,
  urlcolor=blue
}

\lstset{
  basicstyle=\ttfamily\small,
  breaklines=true,
  frame=single,
  columns=fullflexible
}

\title{Эвристические и ML-подходы к решению задачи 3D Bin Packing}
\authorfull{Иртюга Илья Владимирович}
\author{Иртюга И. В.}
\Ministry{Министерство науки и высшего образования Российской Федерации}
\KursWork
\makeatletter
\@jobspecfalse
\makeatother

\faculty{фундаментальных наук}
\facultyshort{ФН}
\chair{математического моделирования}
\chairno{}
\group{ФН12-81Б}

\chief{Фетисов Д. А.}
\FirstConsultant{}
\inspector{}

\workyear{2026}

\setcontentsname{СОДЕРЖАНИЕ}
\setrefsname{Список использованных источников}

\begin{document}
\ifusekprj
  \maketitle
\else
  \begin{titlepage}
    \makeatletter
    \centering
    {\large Курсовая работа\par}
    \vspace{1.5cm}
    {\bfseries\LARGE \@title \par}
    \vspace{1.5cm}
    {\large Автор: \@author\par}
    \vspace{0.5cm}
    {\large Группа: \@group\par}
    \vspace{0.5cm}
    {\large Кафедра: \@chair\par}
    \vspace{0.5cm}
    {\large Руководитель: \@chief\par}
    \vfill
    {\large \@workyear\par}
    \makeatother
  \end{titlepage}
\fi
\tableofcontents
\pagebreak

\section{Введение}

Задача упаковки объектов в ограниченный контейнер относится к классическим комбинаторным задачам оптимизации и известна как bin packing problem~\cite{wiki_bin_packing,martello2000}. В трехмерном варианте требуется разместить набор прямоугольных параллелепипедов внутри контейнера так, чтобы не нарушались геометрические ограничения и одновременно достигалась высокая плотность заполнения. Для логистики и складских систем такая задача имеет прикладное значение: качество укладки влияет на число используемых паллет, устойчивость груза и транспортные издержки~\cite{isa2016}.

В рассматриваемом проекте реализован программный стенд для задачи автоматической укладки коробов на паллету. Репозиторий включает генерацию синтетических задач, набор алгоритмов упаковки, валидатор решений и сохраненные benchmark-результаты~\cite{project_repo}. По содержанию он близок к практическим постановкам соревнований по трехмерной упаковке, где качество решения определяется не только плотностью заполнения, но и временем работы алгоритма~\cite{kaggle_sleigh}.

С вычислительной точки зрения рассматриваемая задача характеризуется сочетанием геометрических, физических и логистических аспектов. В ходе анализа постановки было установлено, что геометрически допустимое размещение должно одновременно удовлетворять требованиям по опоре, массе и правилам обращения с хрупкими товарами. В связи с этим в работе исследуется не только построение допустимой конфигурации, но и получение решения с высокими значениями нескольких критериев качества.

\textbf{Цель работы} состоит в анализе реализованных в проекте эвристических и гибридных методов решения задачи 3D bin packing и в сравнении их качества на имеющихся наборах данных.

\textbf{Для достижения цели решаются следующие задачи:}
\begin{enumerate}
  \item формализовать постановку задачи и метрику качества;
  \item описать структуру репозитория и используемые данные;
  \item рассмотреть реализованные алгоритмы упаковки;
  \item сопоставить методы по benchmark-результатам;
  \item определить особенности текущей реализации и направления дальнейшего развития.
\end{enumerate}

\textbf{Объект исследования} --- программная система для решения задачи 3D bin packing на паллете.

\textbf{Предмет исследования} --- эвристические, метаэвристические, нейросетевые и гибридные методы укладки, а также критерии оценки качества решения.

\textbf{Практическая значимость} работы заключается в том, что рассматриваемый проект формирует воспроизводимую основу для экспериментального сравнения алгоритмов упаковки и может служить базой для учебных исследований и дальнейшего развития прикладных логистических решений.

\pagebreak

\section{Постановка задачи}

\subsection{Формальная модель}

Рассматривается одна паллета с заданными размерами $L \times W \times H$ и ограничением по массе. На вход также подается набор типов коробок; для каждого типа известны длина, ширина, высота, масса одного экземпляра, количество экземпляров и логические признаки ограничений~\cite{project_repo}. Таким образом, каждая задача представляет собой дискретный набор однотипных прямоугольных грузов, которые требуется разместить внутри ограниченного объема.

Для каждого экземпляра коробки алгоритм должен определить:
\begin{itemize}
  \item координаты нижнего ближнего левого угла;
  \item фактические размеры после поворота;
  \item код выбранной ориентации.
\end{itemize}

Если часть коробок разместить не удается, они остаются вне решения. Постановка относится к трехмерной задаче упаковки прямоугольных блоков и является NP-трудной, поэтому на практике применяются эвристики, локальный поиск и гибридные методы~\cite{martello2000,faroe2003}.

\subsection{Ограничения}

В проекте решение проверяется валидатором, который накладывает следующие жесткие ограничения~\cite{project_repo}:
\begin{enumerate}
  \item каждая коробка должна целиком находиться внутри паллеты;
  \item коробки не должны пересекаться по объему;
  \item суммарная масса не должна превышать допустимое значение;
  \item для коробок с признаком \texttt{strict\_upright} запрещены повороты, меняющие исходную высоту;
  \item для коробок, расположенных выше пола паллеты, требуется опора не менее чем на 60\% площади основания;
  \item тяжелые коробки не должны размещаться над хрупкими.
\end{enumerate}

Условие опоры можно записать в виде
\[
\frac{S_{\text{support}}}{S_{\text{base}}} \ge 0.6,
\]
где $S_{\text{base}}$ --- площадь основания коробки, а $S_{\text{support}}$ --- суммарная площадь ее контакта с лежащими ниже коробками. Это ограничение приближенно учитывает физическую устойчивость укладки.

\subsection{Метрика качества}

Итоговая оценка решения задается взвешенной суммой четырех компонент~\cite{project_repo}:
\[
\text{score} =
0.50 \cdot U_{\text{vol}} +
0.30 \cdot C_{\text{items}} +
0.10 \cdot F_{\text{frag}} +
0.10 \cdot T_{\text{time}}.
\]

Здесь:
\begin{itemize}
  \item $U_{\text{vol}}$ --- утилизация объема паллеты;
  \item $C_{\text{items}}$ --- доля размещенных экземпляров;
  \item $F_{\text{frag}}$ --- оценка соблюдения ограничений хрупкости;
  \item $T_{\text{time}}$ --- временной балл.
\end{itemize}

Таким образом, задача является многокритериальной. При оценке решения одновременно учитываются плотность заполнения контейнера, полнота размещения заказа, корректность укладки и соблюдение временного бюджета.

Из приведенной формулы следует важная практическая особенность. Наибольший вклад в итоговую оценку вносят утилизация объема и полнота размещения, при этом временной фактор и соблюдение ограничений хрупкости также существенно влияют на итоговый результат. Следовательно, в рамках исследования сопоставление алгоритмов проводится по совокупности показателей, а не только по плотности укладки.

\pagebreak

\section{Структура проекта и данные}

\subsection{Состав репозитория}

Репозиторий организован как экспериментальная среда для исследования алгоритмов упаковки~\cite{project_repo}. Его основные компоненты:
\begin{itemize}
  \item \texttt{solvers} --- single-pallet решатели;
  \item \texttt{dataset\_generation} --- генерация синтетических задач;
  \item \texttt{datasets} --- готовые наборы входных данных;
  \item \texttt{results} --- сохраненные результаты запусков;
  \item \texttt{multipallet} --- расширение алгоритмов на многопаллетный сценарий;
  \item \texttt{validator.py} --- проверка корректности и вычисление score.
\end{itemize}

Такая структура позволяет исследовать задачу не только как отдельный алгоритм, но и как полный цикл: генерация данных, запуск методов, валидация и сравнение результатов.

Для учебного и исследовательского проекта это имеет принципиальное значение, поскольку наличие единой инфраструктуры обеспечивает воспроизводимость вычислительного эксперимента. При запуске алгоритмов в общем контуре и оценке единым валидатором формируется корректная база для сопоставления характеристик различных подходов.

\subsection{Синтетические данные}

Так как в проекте используются синтетические данные, важную роль играет генератор сценариев~\cite{project_repo}. Он формирует задачи с различными типами товаров и паллет, что позволяет моделировать разные режимы сложности. В наборы входят тяжелые грузы, хрупкие товары, смеси объектов с ограничениями на ориентацию и сценарии с большим количеством коробок.

В каталоге \texttt{datasets} представлены наборы: 
\texttt{dataset\_100.json}, 
\newline
\texttt{dataset\_150.json}, \texttt{dataset\_200\_2.json} и \texttt{dataset\_benchmark.json}. Наиболее интересным для сравнения методов является \texttt{dataset\_benchmark.json}, поскольку он содержит наиболее реалистичные задачи и лучше показывает компромисс между качеством упаковки и временем работы.

Использование нескольких наборов данных имеет существенное значение, поскольку различные сценарии по-разному нагружают алгоритмы. Для одних задач определяющим является размещение тяжелых и габаритных коробок, для других --- корректное формирование верхних слоев, для третьих --- обработка большого количества разнотипных объектов. По результатам эксперимента было получено более надежное представление о поведении решателей на совокупности датасетов, чем при анализе на одном наборе примеров.

\subsection{Контур benchmark-экспериментов}

Для сравнения методов в проекте предусмотрены сценарии массового запуска решателей и последующей агрегации результатов~\cite{project_repo}. По каждому алгоритму сохраняются:
\begin{itemize}
  \item средний итоговый score;
  \item средняя утилизация объема;
  \item средняя полнота размещения;
  \item средний балл за хрупкость;
  \item временная компонента метрики.
\end{itemize}

Наличие общего benchmark-контура позволяет опираться не на отдельные удачные примеры, а на усредненную оценку поведения методов на серии задач.

\pagebreak

\section{Реализованные алгоритмы}

\subsection{Greedy}

Жадный алгоритм служит базовым методом проекта~\cite{project_repo}. Его идея состоит в последовательном размещении коробок по одной с выбором локально допустимой позиции. Для каждой коробки перебираются возможные ориентации и кандидаты на размещение, после чего выбирается вариант, который удовлетворяет ограничениям и располагает объект как можно ниже и компактнее.

По идеологии этот подход близок к классическим эвристикам first fit decreasing и bottom-left-fill, широко применяемым в задачах упаковки~\cite{wiki_bin_packing}. В рассматриваемой реализации жадный метод дополнен инженерными ускорениями, позволяющими быстро проверять пересечения и опору без полного перебора уже размещенных коробок.

\texttt{Greedy} характеризуется высокой скоростью работы и простой интерпретируемой логикой построения решения. Особенностью метода является последовательное принятие локальных решений, определяющих структуру итоговой укладки. В результате данный подход выступает в работе в качестве базового ориентира для дальнейшего сравнения более сложных методов.

\subsection{MaxRects}

Отдельный решатель проекта использует послойную схему упаковки и двумерный алгоритм \textit{MaxRects}, известный по задачам размещения прямоугольников и соревнованиям по упаковке~\cite{kaggle_sleigh,project_repo}. На каждом уровне по высоте выбирается толщина слоя, а затем коробки размещаются в плоскости основания как 2D-прямоугольники. После этого формируется следующий слой.

Такой подход ориентирован на быстрое построение плотной послойной структуры. Его характерной особенностью является малая вычислительная стоимость. Качество решения в значительной степени определяется на этапе выбора структуры слоев, что делает данный метод эффективным для быстрого формирования регулярной укладки.

Проведенный анализ показал, что MaxRects эффективно работает в сценариях, где объекты могут быть рационально распределены по слоям с последующим плотным двумерным заполнением каждого уровня. По результатам расчетов было получено, что данный метод сочетает высокую скорость работы с хорошими показателями плотности укладки.

\subsection{LNS}

Основным методом проекта является \texttt{LNS} (\textit{Large Neighborhood Search})
\newline
~\cite{faroe2003,project_repo}. Он начинается с исходной укладки, после чего многократно выполняет цикл разрушения и восстановления решения. На шаге разрушения часть уже размещенных коробок удаляется, а на шаге восстановления освободившееся пространство заново заполняется.

Ключевая идея метода состоит в том, что на каждой итерации выполняется локальный пересмотр части текущей укладки вместо полного построения решения заново. В проекте восстановление реализовано процедурой с использованием beam search поверх частичного решения. За счет этого LNS формирует более качественные конфигурации по сравнению с однократным жадным построением.

С практической точки зрения LNS занимает промежуточное положение между быстрыми эвристиками и более сложными методами локального поиска. В проведенных экспериментах было показано, что данный алгоритм обеспечивает высокий уровень итоговой метрики при сохранении приемлемого времени вычислений.

Для данного проекта это определяет особую значимость LNS: метод использует качественное начальное приближение и далее последовательно улучшает найденную укладку. Проведенное исследование показало, что повторный анализ частичного решения и перестройка его фрагментов обеспечивают рост плотности и полноты размещения в пределах допустимого времени расчета.

\subsection{GAN+GA}

Гибридный метод \texttt{GAN+GA} сочетает генетический алгоритм и генеративную модель~\cite{ganga_paper,project_repo}. Хромосома кодирует порядок рассмотрения коробок и выбор ориентаций, а качество особи оценивается через получаемое решение упаковки. Генетические операторы позволяют исследовать пространство перестановок, тогда как генеративная часть используется для построения новых перспективных конфигураций.

Такой подход представляет исследовательский интерес, поскольку объединяет эволюционный поиск и идеи машинного обучения. В репозитории реализована гибридная схема, расширяющая набор рассматриваемых методов решения задачи.

С методологической точки зрения наличие такого решателя существенно расширяет экспериментальное пространство проекта. Реализация GAN+GA обеспечивает возможность сопоставления классических инженерных эвристик с современными ML-ориентированными подходами в едином вычислительном контуре.

\pagebreak

\section{Экспериментальное сравнение}

\subsection{Сравнение итогового score}

Для четырех single-pallet алгоритмов в репозитории имеются общие результаты трех наборов данных: \texttt{dataset\_100.json}, \texttt{dataset\_200\_2.json} и \texttt{dataset\_benchmark.json}. Средние значения итогового score приведены в таблице~\ref{tab:score-main}.

\begin{table}[H]
\centering
\caption{Средний итоговый score на общих наборах}
\label{tab:score-main}
\begin{tabular}{lccc}
\toprule
\textbf{Алгоритм} & \textbf{dataset\_100} & \textbf{dataset\_200\_2} & \textbf{dataset\_benchmark} \\
\midrule
Greedy & 0.6842 & 0.7120 & 0.6549 \\
MaxRects & 0.7104 & 0.7117 & 0.7124 \\
GAN+GA & 0.7012 & 0.7327 & 0.7161 \\
LNS & \textbf{0.7111} & \textbf{0.7434} & \textbf{0.7298} \\
\bottomrule
\end{tabular}
\end{table}

Из таблицы видно, что LNS показывает лучший средний score на всех трех общих наборах. Наиболее выраженное преимущество наблюдается на \texttt{dataset\_benchmark.json}, где рассматриваются более крупные и сложные задачи. По результатам эксперимента было получено подтверждение высокой эффективности локальной перестройки укладки при росте числа объектов.

Преимущество LNS носит устойчивый характер и наблюдается на нескольких наборах данных. Наибольшая разница фиксируется в задачах с большим пространством решений, что подтверждает результативность метаэвристического механизма частичной перестройки уже построенной укладки.

\subsection{Анализ компонент метрики}

Более детально различия между методами видны на уровне отдельных компонент score. Для набора \texttt{dataset\_benchmark.json} результаты представлены в таблице~\ref{tab:components-benchmark}.

\begin{table}[H]
\centering
\caption{Средние компоненты метрики на \texttt{dataset\_benchmark.json}}
\label{tab:components-benchmark}
\begin{tabular}{lccccc}
\toprule
\textbf{Алгоритм} & \textbf{Score} & \textbf{Vol.} & \textbf{Coverage} & \textbf{Fragility} & \textbf{Time} \\
\midrule
Greedy & 0.6549 & 0.6155 & 0.4906 & 1.0000 & 1.0000 \\
MaxRects & 0.7124 & 0.7513 & 0.6160 & 0.5198 & 1.0000 \\
GAN+GA & 0.7161 & 0.7156 & 0.5540 & 1.0000 & 0.9205 \\
LNS & \textbf{0.7298} & 0.7172 & 0.6036 & 0.9833 & 0.9175 \\
\bottomrule
\end{tabular}
\end{table}

По результатам сравнительного анализа было установлено несколько характерных особенностей.

\begin{itemize}
  \item \texttt{Greedy} максимально быстр и аккуратно соблюдает ограничения хрупкости, но уступает по утилизации объема и покрытию заказа.
  \item \texttt{MaxRects} достигает очень высокой плотности упаковки, однако заметно теряет в компоненте хрупкости.
  \item \texttt{GAN+GA} сохраняет высокий балл по хрупкости и улучшает объемную утилизацию, но теряет часть времени на более сложный поиск.
  \item \texttt{LNS} не максимизирует ни одну компоненту любой ценой, однако дает лучший суммарный баланс между утилизацией, полнотой размещения и временем.
\end{itemize}

Именно указанный баланс объясняет лидерство LNS по итоговому score.

С прикладной точки зрения полученные результаты позволяют охарактеризовать MaxRects как метод, ориентированный на скорость и геометрическую плотность, Greedy --- как устойчивый базовый решатель, а LNS --- как наиболее универсальный алгоритм по совокупности показателей. Реализация GAN+GA дополняет данное сравнение гибридным ML-ориентированным подходом.

\subsection{Сравнение по времени}

Фактическое время решения также существенно влияет на итоговую метрику. Наименьшее время расчета демонстрируют \texttt{MaxRects} и \texttt{Greedy}. Алгоритмы \texttt{LNS} и \texttt{GAN+GA} требуют большего вычислительного времени, при этом сохраняют высокий временной балл. По результатам эксперимента было установлено, что в данной постановке наибольшую практическую ценность представляет алгоритм, обеспечивающий высокий итоговый score при соблюдении допустимого временного бюджета.

\pagebreak

\section{Многопаллетное расширение}

В проекте реализовано расширение на случай нескольких паллет~\cite{project_repo}. Такая постановка ближе к реальной логистике, поскольку полный заказ нередко требует распределения между несколькими контейнерами. В этом случае алгоритм последовательно использует новые паллеты до полного размещения допустимых коробок.

Содержательно многопаллетный режим развивается из single-pallet схемы. Жадная версия повторно запускает базовый алгоритм на остатке коробок, а LNS переносит на каждую паллету ту же идею локального улучшения. В результате в репозитории сформирован отдельный контур решения многопаллетной постановки.

В рамках данной работы многопаллетное расширение рассматривается как самостоятельное направление развития проекта, тогда как центральным объектом анализа выступает single-pallet постановка.

Наличие многопаллетной ветки имеет существенное значение для оценки зрелости проекта. Реализованное расширение демонстрирует переход от локальной задачи плотной укладки к более общей задаче распределения груза между несколькими паллетами.

\pagebreak

\section{Особенности реализации и направления развития}

\subsection{Особенности текущей реализации}

Текущая версия проекта характеризуется рядом существенных особенностей, определяющих направления дальнейшего исследования.

\begin{enumerate}
  \item benchmark-эксперименты организованы на нескольких наборах данных с акцентом на сопоставимые сценарии запуска алгоритмов;
  \item проверка устойчивости основана на критерии опоры не менее 60\% площади основания;
  \item гибридные стохастические методы реализованы с настраиваемыми параметрами поиска;
  \item многопаллетный режим представлен как отдельное направление вычислительного эксперимента.
\end{enumerate}

\subsection{Перспективы развития}

Наиболее естественные направления дальнейшей работы:
\begin{enumerate}
  \item усложнение физической модели устойчивости;
  \item расширение метрики качества с учетом числа используемых паллет;
  \item статистический анализ устойчивости стохастических методов;
  \item развитие более сильных destroy/repair-операторов в LNS;
  \item исследование новых гибридных схем, совмещающих эвристику и обучение.
\end{enumerate}

Особенно значимым направлением является развитие LNS, поскольку уже в текущей реализации именно этот метод обеспечивает наилучший компромисс между качеством и временем.

Кроме того, дальнейшее развитие проекта связано не только с усложнением отдельных решателей, но и с совершенствованием экспериментальной методологии. Существенный интерес представляет систематический анализ чувствительности алгоритмов к типу датасета, числу объектов, средней массе коробок и доле хрупких товаров. Такой анализ обеспечивает более точное понимание режимов наибольшей эффективности каждого метода и позволяет определить наиболее результативные направления инженерного развития.

\pagebreak

\section{Заключение}

В работе была рассмотрена задача 3D bin packing для укладки коробок на паллету и проанализирован репозиторий, содержащий набор алгоритмов, генерацию данных, валидатор и benchmark-результаты. В ходе исследования было показано, что рассматриваемая постановка выходит за рамки простой геометрической упаковки, поскольку учитывает массу, ограничения ориентации, устойчивость и хрупкость груза.

Сравнение реализованных методов показало различие их функциональных ролей. \texttt{Greedy} выступает быстрым базовым решением, \texttt{MaxRects} обеспечивает очень высокую скорость и плотную послойную укладку, \texttt{GAN+GA} представляет гибридный исследовательский подход, а \texttt{LNS} демонстрирует наилучший общий баланс между утилизацией объема, полнотой размещения и временем работы.

По benchmark-результатам было установлено, что именно LNS является сильнейшим single-pallet методом проекта: он показывает лучший средний итоговый score на всех общих наборах данных и особенно уверенно проявляет себя на наиболее сложных задачах. Полученные результаты позволяют рассматривать LNS как основной практически значимый алгоритм текущей реализации, а остальные методы --- как важные точки сравнения и направления дальнейшего исследования.

\appendix
\section{Сводная таблица benchmark-результатов}

\small
\begin{longtable}{@{}p{0.16\textwidth}p{0.14\textwidth}@{\hspace{1.0em}}ccccc@{}}
\toprule
\textbf{Алгоритм} & \textbf{Набор} & \textbf{Score} & \textbf{Vol.} & \textbf{Coverage} & \textbf{Fragility} & \textbf{Time} \\
\midrule
\endfirsthead
\toprule
\textbf{Алгоритм} & \textbf{Набор} & \textbf{Score} & \textbf{Vol.} & \textbf{Coverage} & \textbf{Fragility} & \textbf{Time} \\
\midrule
\endhead
Greedy & dataset100 & 0.6842 & 0.5961 & 0.6204 & 1.0000 & 1.0000 \\
Greedy & dataset200 & 0.7120 & 0.6228 & 0.6686 & 1.0000 & 1.0000 \\
Greedy & datasetBench & 0.6549 & 0.6155 & 0.4906 & 1.0000 & 1.0000 \\
MaxRects & dataset100 & 0.7104 & 0.7005 & 0.6825 & 0.5545 & 1.0000 \\
MaxRects & dataset200 & 0.7117 & 0.7000 & 0.6809 & 0.5737 & 1.0000 \\
MaxRects & datasetBench & 0.7124 & 0.7513 & 0.6160 & 0.5198 & 1.0000 \\
GAN+GA & dataset100 & 0.7012 & 0.6278 & 0.6271 & 1.0000 & 0.9910 \\
GAN+GA & dataset150 & 0.7262 & 0.6571 & 0.6589 & 1.0000 & 1.0000 \\
GAN+GA & dataset200 & 0.7327 & 0.6700 & 0.6591 & 1.0000 & 1.0000 \\
GAN+GA & datasetBench & 0.7161 & 0.7156 & 0.5540 & 1.0000 & 0.9205 \\
LNS & dataset100 & 0.7111 & 0.6384 & 0.6580 & 0.9600 & 0.9850 \\
LNS & dataset200 & 0.7434 & 0.6712 & 0.7055 & 0.9747 & 0.9865 \\
LNS & datasetBench & 0.7298 & 0.7172 & 0.6036 & 0.9833 & 0.9175 \\
\bottomrule
\end{longtable}
\normalsize

\pagebreak

\begin{thebibliography}{99}

\bibitem{wiki_bin_packing}
Bin packing problem // Wikipedia: The Free Encyclopedia [Электронный ресурс]. URL: \url{https://en.wikipedia.org/wiki/Bin_packing_problem} (дата обращения: 27.04.2026).

\bibitem{isa2016}
Статья из журнала Института системного анализа РАН [Электронный ресурс]. URL: \url{http://www.isa.ru/aidt/2016-03/31_43.pdf} (дата обращения: 27.04.2026).

\bibitem{kaggle_sleigh}
Packing Santa's Sleigh // Kaggle [Электронный ресурс]: соревнование по задаче трехмерной упаковки. URL: \url{https://www.kaggle.com/competitions/packing-santas-sleigh/overview} (дата обращения: 27.04.2026).

\bibitem{project_repo}
DrogoZZZZ. \textit{3D Bin Packing} [Электронный ресурс]: GitHub-репозиторий проекта. URL: \url{https://github.com/DrogoZZZZ/3D_bin_packing} (дата обращения: 27.04.2026).

\bibitem{martello2000}
Martello S., Pisinger D., Vigo D. The Three-Dimensional Bin Packing Problem // \textit{Operations Research}. 2000. Vol. 48. No. 2. P. 256--267. DOI: \url{https://doi.org/10.1287/opre.48.2.256.12386}.

\bibitem{faroe2003}
Faroe O., Pisinger D., Zachariasen M. Guided Local Search for the Three-Dimensional Bin-Packing Problem // \textit{INFORMS Journal on Computing}. 2003. Vol. 15. No. 3. P. 267--283.

\bibitem{ganga_paper}
Zhang Y., Wang X., Liu Z. A Novel Approach for Solving 3D Bin Packing Problem by Integrating a Generative Adversarial Network with a Genetic Algorithm // \textit{Scientific Reports}. 2024. DOI: \url{https://doi.org/10.1038/s41598-024-56699-7}.

\end{thebibliography}

\end{document}
